%!TEX program = xelatex
\documentclass[a4paper,12pt]{article}
\usepackage{fontspec}\defaultfontfeatures{Ligatures=TeX}
\usepackage[a4paper,vmargin={3cm,3.5cm},hmargin={3cm,3cm}]{geometry}
%-----------------------------------------------------------------------------%
\usepackage{settings}
%-----------------------------------------------------------------------------%
%%% Title %%%
\title{G6411 Recitation Notes: \glsentrylong{mds}}
\author{Jesse Chieh Chen\\\href{mailto:cc5229@columbia.edu}{cc5229@columbia.edu}}
\date{\today}
%-----------------------------------------------------------------------------%

\begin{document}

\appendix
\section{\glsentrylong{mds}}

Time series observations need not be independent for their averages to satisfy a central limit theorem.
A useful alternative is that each new observation has conditional mean zero given the past.
This condition will be especially useful for the products of regressors and errors that appear in estimation.

\subsection{Unpredictable in Mean}

Let $\mathcal F_t$ denote the information available through date $t$.
Information accumulates over time: $\mathcal F_{t-1}\subseteq\mathcal F_t$.
For example, $\mathcal F_t$ may contain all observations and shocks up to date $t$.

\begin{definition}[\gls{mds}]
	A sequence $\{g_t\}$, with $g_t\in\reals^k$, is an \gls{mds} relative to $\{\mathcal F_t\}$ if
	$g_t$ is known at date $t$, $\E\|g_t\|<\infty$, and
	\begin{align}\label{eq-mds}
		\E(g_t\given\mathcal F_{t-1})=0.
	\end{align}
\end{definition}

\begin{remark}[Interpretation]
	Even using all the information in $\mathcal F_{t-1}$,
	the best mean prediction of $g_t$ is zero.
\end{remark}

\begin{remark}[Past Information]
	The information set matters:
	a condition given the history of $g_t$ alone need not remain true after adding other variables to that history.
\end{remark}

\begin{remark}[Where is the ``Difference''?]
	The name comes from the partial sum $S_t=\sum_{s=1}^t g_s$.
	Then, we have $\E(S_t\given\mathcal F_{t-1})=S_{t-1}$, so $S_t$ is a \emph{martingale} and $g_t=S_t-S_{t-1}$ is its difference.
\end{remark}

\begin{proposition}[Mean Zero and No Serial Correlation]\label{prop-mds-orthogonal}
	An \gls{mds} has $\E g_t=0$.
	If it has finite second moments, then for $s<t$,
	\begin{align}
		\E(g_tg_s\T)=0.
	\end{align}
\end{proposition}
\begin{proof}
	The \gls{lie} gives $\E g_t=\E[\E(g_t\given\mathcal F_{t-1})]=0$.
	Since $g_s$ is known at date $t-1$,
	\begin{align*}
		\E(g_tg_s\T)
		&=\E\left[\E(g_t\given\mathcal F_{t-1})g_s\T\right]=0.
		\qedhere
	\end{align*}
\end{proof}

\begin{remark}
	More generally, $g_t$ is orthogonal to any function of past information for which the product is integrable.
	This is stronger than being uncorrelated with individual past observations.
	An \gls{mds} need not be stationary or have constant variance;
	a scalar \gls{mds} with constant finite positive variance is white noise.
\end{remark}

\begin{example}[Dependence and Conditional Heteroskedasticity]
	Let $z_t\iidto\normaldistribution(0,1)$ and set $e_t=z_{t-1}z_t$.
	With $\mathcal F_t$ consisting of the history of $z_t$ through date $t$,
	\begin{align}
		\E(e_t\given\mathcal F_{t-1})=0,
		\qquad \E(e_t^2\given\mathcal F_{t-1})=z_{t-1}^2.
	\end{align}
	Thus $e_t$ is an \gls{mds} with variance one but a time-varying conditional variance.
	Dependence remains: $\E(e_t^2e_{t-1}^2)=3\neq1=\E(e_t^2)\E(e_{t-1}^2)$.
\end{example}

\begin{example}[White Noise Need Not Be an \gls{mds}]
	With the same independent standard normal variables, let $u_t=z_t+z_{t-1}^2-1$.
	Then $\E u_t=0$, $\var(u_t)=3$, and all nonzero-lag covariances vanish:
	at lag one, the only potentially nonzero term is $\E[(z_{t-1}^2-1)z_{t-1}]=0$.
	Nevertheless,
	\begin{align}
		\E(u_tu_{t-1}^2)=\E[(z_{t-1}^2-1)z_{t-1}^2]=2\neq0.
	\end{align}
	So $u_t$ is white noise but not an \gls{mds}, even relative to its own past.
	A nonlinear function of the past can help predict it.
\end{example}

\subsection{The Central Limit Theorem}

The following is a convenient stationary version of the \gls{mds} \gls{clt}.
Ergodicity ensures that time averages recover population moments;
it is an additional condition, not a consequence of being an \gls{mds}.

\begingroup
\tcbset{lines before break=4}
\begin{theorem}[\gls{mds} \gls{clt}]\label{thm-mds-clt}
	Let $\{g_t\}$ be a strictly stationary and ergodic \gls{mds} in $\reals^k$.
	Suppose $\E\|g_t\|^2<\infty$, and define $\Omega=\E(g_tg_t\T)$.
	Then, as $T\to\infty$,
	\begin{align}\label{eq-mds-clt}
		\frac1{\sqrt T}\sum_{t=1}^Tg_t
		=\sqrt T\,\overline g_T
		\dto\normaldistribution(0,\Omega).
	\end{align}
	For a scalar sequence, the asymptotic variance is simply $\E g_t^2$.
\end{theorem}
\endgroup

\begin{remark}[Where Is the \gls{lrv}?]
	There is no missing covariance adjustment.
	Writing $\Gamma(j)=\E(g_tg_{t-j}\T)$, \autoref{prop-mds-orthogonal} gives
	\begin{align}
		\sum_{j=-\infty}^{\infty}\Gamma(j)=\Gamma(0)=\Omega.
	\end{align}
	Thus the \gls{lrv} equals the contemporaneous variance for an \gls{mds}.
	Zero serial covariances explain this variance formula, but do not by themselves establish a \gls{clt}.
\end{remark}

\begin{remark}[Estimating the Variance]
	If $g_t$ is observed, ergodicity and the finite second moment give
	\begin{align}
		\widehat\Omega=\frac1T\sum_{t=1}^Tg_tg_t\T\pto\Omega.
	\end{align}
	No lagged products are needed.
	In the scalar case with $\Omega>0$, Slutsky's theorem gives
	$\widehat\Omega^{-1/2}\sqrt T\,\overline g_T\dto\normaldistribution(0,1)$.
\end{remark}

\subsection{Why This Helps with Regression}

Consider $Y_t=X_t\T\beta+e_t$, with $X_t\in\reals^k$.
Suppose $X_t$ is known at date $t-1$ and $\E(e_t\given\mathcal F_{t-1})=0$.
Provided $\E\|X_te_t\|<\infty$, the vector $g_t=X_te_t$ satisfies
\begin{align}\label{eq-score-mds}
	\E(g_t\given\mathcal F_{t-1})
	=X_t\E(e_t\given\mathcal F_{t-1})=0.
\end{align}
Hence multiplying an \gls{mds} error by a regressor known from the past produces an \gls{mds} score.
If $\{X_te_t\}$ is also strictly stationary and ergodic, with $\E(e_t^2\|X_t\|^2)<\infty$,
\autoref{thm-mds-clt} gives
\begin{align}
	\frac1{\sqrt T}\sum_{t=1}^T g_t\dto\normaldistribution(0,\Omega),
	\qquad \Omega=\E g_t g_t\T.
\end{align}
Let $\widehat\beta$ be the \gls{ols} estimator.
If $\widehat Q_{xx}=T^{-1}\sum_tX_tX_t\T\pto Q_{xx}$ with $Q_{xx}$ positive definite, then
\begin{align}
	\sqrt T(\widehat\beta-\beta)
	=\widehat Q_{xx}\inv\frac1{\sqrt T}\sum_{t=1}^T g_t
	\dto\normaldistribution(0,V),\qquad
	V=Q_{xx}\inv\Omega Q_{xx}\inv.
\end{align}

\begin{remark}[Application to ADL Models]
	For an \gls{adl} model of order $(p,q)$, take
	\begin{align*}
		X_t=(1,Y_{t-1},\ldots,Y_{t-p},Z_t,Z_{t-1},\ldots,Z_{t-q})\T,
	\end{align*}
	and let $\mathcal F_t$ be the history of $Y_t$ and $Z_t$ through date $t$.
	Since $Z_t$ need not be known at date $t-1$, assume
	$\E(e_t\given\mathcal F_{t-1},Z_t)=0$.
	Now $X_t$ is known given $(\mathcal F_{t-1},Z_t)$, so the \gls{lie} gives
	\begin{align*}
		\E(g_t\given\mathcal F_{t-1})
		&=\E\left[X_t\E(e_t\given\mathcal F_{t-1},Z_t)\given\mathcal F_{t-1}\right]=0.
	\end{align*}
	Thus the ADL score $g_t=X_te_t$ is an \gls{mds} relative to $\{\mathcal F_t\}$.
\end{remark}

\begin{remark}[Conditional Homoskedasticity]
	Let $\sigma_t^2=\E(e_t^2\given\mathcal F_{t-1})$.
	Because $X_t$ is known from the past,
	\begin{align}
		\Omega=\E g_tg_t\T = \E X_tX_t\T\sigma_t^2.
	\end{align}
	For ADL models with contemporaneous $Z_t$, use $\sigma_t^2=\E(e_t^2\given\mathcal F_{t-1},Z_t)$ instead;
	the same calculation applies.
	If $\sigma_t^2=\sigma^2$, then $\Omega=\sigma^2Q_{xx}$ and $V=\sigma^2Q_{xx}\inv$.
	Under conditional heteroskedasticity, the \gls{mds} \gls{clt} still applies,
	but we retain the sandwich variance $V$.
\end{remark}

\vfill
\section*{Acronyms}
\printglossaries

\end{document}
